% ============================================================
% Table: Violation Sequencing — What Follows a Mining vs Non-mining MR Violation?
% Unit: MR violation (unique VIOLATION_ID). Next violation = highest-severity
% violation at the next distinct begin_date for the same PWSID.
% N: mining MR = 1,273
% ============================================================

\begin{table}[htbp]
\centering
\caption{Violation Following an Arsenic, Nitrate, or Inorganic Chemical MR Violation, Downstream CWSs, 1985--2005}
\label{tab:viol_sequence}
\small
\begin{tabular}{lr}
\hline\hline
\textbf{Next violation} & \textbf{Mining MR} \\
 & \textit{(N=1,273)} \\
\hline
\addlinespace[4pt]
\multicolumn{2}{l}{\textit{No subsequent violation}} \\
\addlinespace[2pt]
No further violation in 1985--2005 (\%) & 2.2 \\
\addlinespace[4pt]
\multicolumn{2}{l}{\textit{Monitoring/reporting violation follows}} \\
\addlinespace[2pt]
MR, same contaminant$^a$ (\%) & 4.9 \\
MR, same contaminant group$^b$ (\%) & 9.8 \\
MR, other contaminant (\%) & 63.1 \\
\addlinespace[4pt]
\multicolumn{2}{l}{\textit{Limit exceedance violation follows}} \\
\addlinespace[2pt]
MCL, same contaminant$^a$ (\%) & 0.2 \\
MCL, same contaminant group$^b$ (\%) & 0.3 \\
MCL, other contaminant (\%) & 6.0 \\
\addlinespace[4pt]
\multicolumn{2}{l}{\textit{Other violation type follows}} \\
\addlinespace[2pt]
Treatment technique (TT) (\%) & 2.1 \\
Other category (\%) & 11.4 \\
\addlinespace[4pt]
\multicolumn{2}{l}{\textit{Time to next violation}} \\
\addlinespace[2pt]
Median days to next violation & 365.0 \\
Mean days to next violation & 674.7 \\
\hline\hline
\end{tabular}
\begin{minipage}{\linewidth}
\vspace{4pt}
\footnotesize
\textit{Notes:} Unit of observation is the MR violation (unique VIOLATION\_ID). Sample: CWS in the coal mining analysis panel whose intake HUC is one step downstream of a mine HUC (downstream-only sample), 1985--2005. Mining-related rules: nitrate (331), arsenic (332), inorganic chemicals (333). Non-mining rules: total coliform (110, 111), surface/groundwater rule (121--123, 140), VOCs (310), SOCs (320). ``Next violation'' is defined as the highest-severity violation beginning on the next distinct begin date for the same PWSID; severity order is MCL mining $>$ MCL non-mining $>$ MCL other $>$ TT $>$ MR mining $>$ MR non-mining $>$ other. When multiple violations share the next date, the highest-severity is selected; this is conservative for detecting MCL escalation (the MCL is surfaced if present). ``Days to next violation'' is missing for violations with no subsequent violation. $^a$Same contaminant = same SDWIS RULE\_CODE (e.g., an arsenic MR followed by an arsenic MCL). $^b$Same contaminant group (but different rule code): e.g., a nitrate MR followed by an arsenic MCL; or a total-coliform MR followed by a surface-water rule MCL. Source: SDWA\_VIOLATIONS\_ENFORCEMENT.csv.
\end{minipage}
\end{table}
